\SourceLeaf{611}\begin{center}[4.]\end{center}
Von der \textit{Geometria Situs}, die \textsc{Leibnitz} erahnte und in die nur zwei Geometer (\textsc{Euler} und \textsc{Vandermonde}) einen schwachen Einblick gewinnen durften, wissen wir nach anderthalb Jahrhunderten noch immer fast nichts und haben auch kaum etwas davon gewonnen.

Eine Hauptaufgabe im \textit{Grenzgebiet} zwischen \textit{Geometria Situs} und \textit{Geometria Magnitudinis} wird darin bestehen, die Umschlingungen zweier geschlossener oder unendlicher Linien zu zählen.

Seien $x,y,z$ die Koordinaten eines unbestimmten Punkts der ersten Linie und $x',y',z'$ die eines solchen Punkts der zweiten, und sei
\[
\int\int\frac{\begin{gathered}
(x'-x)(\mathrm dy\,\mathrm dz'-\mathrm dz\,\mathrm dy')\\
+(y'-y)(\mathrm dz\,\mathrm dx'-\mathrm dx\,\mathrm dz')\\
+(z-z')(\mathrm dx\,\mathrm dy'-\mathrm dy\,\mathrm dx')
\end{gathered}}{[(x'-x)^2+(y'-y)^2+(z'-z)^2]^{\frac{3}{2}}}=V
\]
Dann ergibt dieses Integral, über beide Linien erstreckt, den Wert
\[
=4m\pi
\]
wobei $m$ die Anzahl der Umschlingungen ist.

Der Wert ist wechselseitig, das heißt, er bleibt beim Vertauschen der beiden Linien gleich.\quad 1833.\ Jan.\ 22.

\begin{center}[5.]\par\medskip\textit{Gesetz des galvanischen Stroms:}\end{center}
1.\quad Positiv ist der Strom, der an der Kontaktstelle mit dem Wasser in der Richtung Zink, Wasser, Kupfer fließt.

2.\quad Sei $RR'$ ein Stromelement, in dem der positive Strom von $R$ nach $R'$ fließt, und sei $P$ ein Punkt, in dem sich ein Element positiven nördlichen magnetischen Fluids befindet.

Das Stromelement $RR'=\mu$ übt dann auf $P$ eine Kraft von der Stärke
\[
=\frac{\mu\,.\,\sin R'RP}{(RR')^2}
\]
aus. Ihre Richtung $PQ$ steht senkrecht auf der Ebene durch $P,R,R'$ und weist nach unten, sofern $P$ rechts von $RR'$ oder $R'$ rechts von $PR$ liegt.
\par
\SourceLeaf{612}
3.\quad Der Wirkung eines geschlossenen Stroms $RR'R''R$ kann man auf folgende Weise eine magnetische Wirkung gleichsetzen. Es begrenze $RR'R''R$ eine beliebige Fläche, auf der man sich nördlichen Magnetismus in beliebiger Verteilung mit der Dichte $=\delta$ ausgebreitet denkt. An jeder Stelle der Fläche errichte man eine unendlich kleine Normale, die direkt proportional zur Intensität des Stroms und umgekehrt proportional zu $\delta$ ist. Sie sei nach oben oder nach unten gerichtet, je nachdem der Strom die Fläche beim Umlauf rechts oder links hat. Die Endpunkte dieser Normalen liegen auf einer zweiten Fläche. Auf dieser Fläche und ihren Teilen breite man genau so viel südlichen Magnetismus aus, wie sich auf der anderen Fläche und ihren entsprechenden Teilen befindet. Diese beiden magnetischen Flächen sind für jeden außerhalb von ihnen liegenden Punkt jenem galvanischen Strom äquivalent.

